\documentclass[11pt,a4paper]{report}
\usepackage[utf8]{inputenc}
\usepackage[T1]{fontenc}
\usepackage{lmodern}
\usepackage[margin=2.3cm]{geometry}
\usepackage{amsmath,amssymb,bm}
\usepackage{booktabs,array,multirow,tabularx}
\usepackage{siunitx}
\usepackage{xcolor}
\usepackage{hyperref}
\usepackage{xurl}
\usepackage{enumitem}
\usepackage{titlesec}
\titleformat{\chapter}[hang]{\Large\bfseries}{\thechapter}{1em}{}
\titlespacing*{\chapter}{0pt}{0pt}{1.2em}
\hypersetup{colorlinks=true,linkcolor=blue!50!black,urlcolor=blue!50!black,citecolor=blue!50!black}
\sisetup{group-separator={\,},group-minimum-digits=5}
\setlist{nosep,leftmargin=1.5em}

\newcommand{\bu}{\bm{u}}
\newcommand{\K}{\mathbf{K}}
\newcommand{\KT}{\mathbf{K}_T}
\newcommand{\KD}{\mathbf{K}_\Delta}
\newcommand{\refval}[1]{\textbf{#1}}

\title{\textbf{Verification Test Cases for the\\Linear Buckling Analysis in 4C}\\[0.5em]
\large Three-dimensional solid elements}
\author{Linear buckling feature, branch \texttt{dev-linear-buckling} of \texttt{m-frey/4C}}
\date{25 September 2026}

\begin{document}
\maketitle
\tableofcontents

%======================================================================
\chapter*{Overview}
\addcontentsline{toc}{chapter}{Overview}
%======================================================================

This document defines four test cases for verifying the linear buckling
analysis of 3D solid elements in 4C. For each case it gives the setup,
the derivation of the reference solution, and the accepted deviation.
All numbers to be matched are collected in Chapter~\ref{ch:results}.

\begin{center}
\begin{tabularx}{\textwidth}{@{}l>{\raggedright}X>{\raggedright}X>{\raggedright\arraybackslash}X@{}}
\toprule
Case & Structure & Reference & Tests mainly \\
\midrule
A & Simply supported plate, uniaxial compression & analytical (Kirchhoff, Mindlin) & plate modes, several $(m,n)$, Neumann load \\
B & Cantilever column, rectangular section & analytical (Euler, Timoshenko) & mode ordering, preload path \\
C & 3D beam-like solid, clamped--pinned & published benchmark (FEniCSx) & code-to-code, identical mesh \\
D & Periodic square-lattice cell & plausibility + exact identities & periodic BCs, Dirichlet load \\
\bottomrule
\end{tabularx}
\end{center}

\paragraph{Eigenvalue problem and notation.}
All cases solve
\begin{equation}
  \bigl(\KT(\bu_0) + \lambda\,\KD\bigr)\,\bm{\phi} = \bm{0},
  \label{eq:evp}
\end{equation}
where $\bu_0$ is the base state ($\bu_0=\bm{0}$ unless stated otherwise),
$\KD$ is the differential (initial-stress) stiffness due to a reference load
pattern $P_\text{ref}$, and the critical load is $P_\text{cr} = P_\text{pre} +
\lambda P_\text{ref}$. Positive $\lambda$ means the load acts as given; negative
$\lambda$ means it acts in reverse.

\paragraph{Differential stiffness.}
4C uses the classical geometric stiffness $\KD = \K_\sigma(\Delta\bm{S})$ of
the perturbation stress only, as Abaqus \texttt{*BUCKLE} and the FEniCSx
benchmark (Case~C) do (design choice D1 of the implementation plan).
A finite-difference tangent,
$\KD \approx [\KT(\bu_0+\varepsilon\Delta\bu)-\KT(\bu_0)]/\varepsilon$, would
additionally contain the initial-displacement term (linear in
$\nabla\Delta\bu$). Its relative effect on $\lambda_i$ is of the order of the
pre-buckling strain at the critical point,
$\varepsilon_{\text{pre},i}\approx\sigma_{\text{cr},i}/E$. This is why that
variant was rejected. The tolerances below apply to the classical
$\KD$.

\paragraph{Mode identification and comparison.}
\begin{itemize}
\item Modes are identified by their dominant displacement component and by
  counting half-waves along the edges or axes.
\item The shape is compared with the analytical mode $\bm{\phi}^\text{an}$ at
  the nodes of the mid-surface or neutral axis, using the modal assurance
  criterion
  $\text{MAC} = (\bm{\phi}^\text{T}\bm{\phi}^\text{an})^2 /
  \bigl((\bm{\phi}^\text{T}\bm{\phi})(\bm{\phi}^{\text{an}\,\text{T}}\bm{\phi}^\text{an})\bigr)$.
\item For degenerate eigenvalues (multiplicity $>1$), only the eigenvalue is
  compared, together with the fact that the mode lies in the analytical
  eigenspace.
\end{itemize}

\paragraph{Units.}
Cases A, B and D use \si{mm}, \si{N} and \si{MPa}. Case C uses the
dimensionless units of the published benchmark.

%======================================================================
\chapter{Case A -- Simply supported plate under uniaxial compression}
\label{ch:plate}
%======================================================================

\section{Problem definition}

A flat rectangular plate of length $a$ (in $x$), width $b$ (in $y$) and
thickness $t$ (in $z$) is simply supported on all four edges and loaded by a
uniform compressive membrane force $N_x$ per unit width on the edges $x=0$
and $x=a$. The plate is modelled as a 3D solid; no periodic boundary
conditions are used.

\begin{center}
\begin{tabular}{@{}lll@{}}
\toprule
Quantity & A1 (square) & A2 (long) \\
\midrule
Length $a$ & \SI{100}{mm} & \SI{300}{mm} \\
Width $b$ & \SI{100}{mm} & \SI{100}{mm} \\
Thickness $t$ & \SI{1}{mm} & \SI{1}{mm} \\
Young's modulus $E$ & \multicolumn{2}{l}{\SI{210000}{MPa}} \\
Poisson's ratio $\nu$ & \multicolumn{2}{l}{0.3} \\
Material & \multicolumn{2}{l}{St.\ Venant--Kirchhoff, \texttt{KINEM nonlinear}} \\
Reference load $N_\text{ref}$ & \multicolumn{2}{l}{\SI{1}{N/mm}, i.e.\ traction $N_\text{ref}/t = \SI{1}{MPa}$} \\
\bottomrule
\end{tabular}
\end{center}

Configuration A1 has well-separated eigenvalues and is the primary test.
Configuration A2 has several closely spaced eigenvalues with different
half-wave numbers $m$. It tests that the eigensolver finds the correct
eigenvalues and returns them in the correct order.

\section{3D model and boundary conditions}

A simple support in the 3D solid has to act on the mid-surface
$z = t/2$ only. Fixing $w$ on the whole edge face would also restrain the
rotation, i.e.\ partially clamp the edge.

\begin{enumerate}
\item \textbf{Out-of-plane support:} $u_z = 0$ on the four mid-surface edge
  lines ($z=t/2$ at $x=0$, $x=a$, $y=0$, $y=b$). This requires a node layer at
  $z=t/2$, i.e.\ an even number of element layers, or quadratic elements with
  mid-thickness nodes.
\item \textbf{Loading:} uniform dead surface traction $p = N_\text{ref}/t$ on
  both loaded faces $x=0$ and $x=a$, acting inward (compression). The load is
  self-equilibrated.
\item \textbf{Remaining in-plane rigid-body modes (statically determinate):}
  \begin{itemize}
  \item $u_x = 0$ and $u_y=0$ at point $P_1=(0,0,t/2)$;
  \item $u_y = 0$ at point $P_2=(a,0,t/2)$.
  \end{itemize}
  The pre-buckling state is then exactly
  $\sigma_{xx} = -N_x/t$, $\sigma_{yy}=\sigma_{xy}=0$. The unloaded edges are
  free in their plane and Poisson expansion is unrestrained.
\item \textbf{Elements:} \texttt{SOLID HEX20} or \texttt{HEX27} with at least
  2 layers through the thickness, or \texttt{HEX8} with EAS
  (\texttt{eas\_full}) and at least 2 layers. Plain \texttt{HEX8} locks
  severely at $b/t=100$ and must not be used.
\item \textbf{Recommended in-plane mesh:} at least 8 quadratic elements per
  half-wave. That is $20\times20$ for A1 and $60\times20$ for A2 (the
  highest listed mode of A2 has $m=8$).
\end{enumerate}

\section{Analytical solution (Kirchhoff plate theory)}

The bending stiffness is $D = E t^3/[12(1-\nu^2)]$. Linearized about the
membrane state $N_x$ (compression positive), the governing equation is
\begin{equation}
  D\,\nabla^4 w + N_x\,\frac{\partial^2 w}{\partial x^2} = 0 .
\end{equation}
The simply supported boundary conditions ($w=0$ and $M_n=0$ on all edges) are
satisfied exactly by the Navier ansatz
$w_{mn} = \sin(m\pi x/a)\,\sin(n\pi y/b)$ with $m,n\in\mathbb{N}$.
Substituting it gives
\begin{equation}
  N_{x,mn} = \frac{\pi^2 D}{b^2}\,k_{mn},
  \qquad
  k_{mn} = \left(\frac{m b}{a} + \frac{n^2 a}{m b}\right)^2 .
  \label{eq:plate}
\end{equation}
The eigenvalues are $\lambda_{mn}=N_{x,mn}/N_\text{ref}$, and $w_{mn}$ is the
corresponding buckling mode, i.e.\ the out-of-plane displacement of the
mid-surface. For $a/b=1$ the minimum is $k_{11}=4$. For $a/b=3$ it is
$k_{31}=4$, with the neighbours $k_{41}=4.340$, $k_{21}=4.694$, $k_{51}=5.138$.

\paragraph{Transverse-shear correction (Mindlin plate theory).}
A 3D solid includes transverse-shear flexibility. For a simply supported
Mindlin plate, the exact relation to the Kirchhoff value is
\begin{equation}
  N^{\text{M}}_{x,mn} = \frac{N_{x,mn}}{1 + D\,\kappa_{mn}^2/(\kappa G t)},
  \qquad \kappa_{mn}^2 = \Bigl(\frac{m\pi}{a}\Bigr)^2+\Bigl(\frac{n\pi}{b}\Bigr)^2,
  \quad G=\frac{E}{2(1+\nu)},\ \kappa=\tfrac56 .
\end{equation}
For $b/t=100$ the correction is below \SI{0.6}{\percent} for all listed
modes. The Mindlin value is the primary reference for the 3D model.

\section{Expected deviations of the 3D model}
\begin{itemize}
\item Remaining 3D effects (thickness change, edge zones of the line support)
  are of order $(t/b)^2$ and negligible.
\item The discretization error of conforming displacement elements usually
  overestimates $\lambda$. It decreases under mesh refinement and dominates
  the total deviation.
\item Pre-buckling strain: $\varepsilon_\text{pre}\approx 75.9/210000
  \approx 3.6\times10^{-4}$ (mode 1), so geometric nonlinearity of the
  pre-buckling state is negligible.
\end{itemize}

%======================================================================
\chapter{Case B -- Cantilever column with rectangular cross-section}
\label{ch:column}
%======================================================================

\section{Problem definition}

A straight prismatic bar of length $L$ with rectangular cross-section
$b\times h$ ($b$ in $y$, $h$ in $z$, $h<b$) is clamped at $x=0$ and loaded at
$x=L$ by a dead axial compressive force $P$.

\begin{center}
\begin{tabular}{@{}ll@{}}
\toprule
Length $L$ & \SI{100}{mm} \\
Cross-section $b\times h$ & \SI{3}{mm} $\times$ \SI{2}{mm}, $A=\SI{6}{mm^2}$ \\
$E$, $\nu$ & \SI{210000}{MPa}, 0.3 \\
Material & St.\ Venant--Kirchhoff, \texttt{KINEM nonlinear} \\
Reference load $P_\text{ref}$ & \SI{1}{N}, i.e.\ traction $1/6$~\si{MPa} on the end face \\
\bottomrule
\end{tabular}
\end{center}

\section{3D model and boundary conditions}
\begin{itemize}
\item \textbf{Clamped end:} $u_x=u_y=u_z=0$ on all nodes of the face $x=0$.
  This also blocks Poisson contraction there, a local effect of order
  $h/L$ on the stress field that is negligible for $\lambda$ at $L/h=50$.
\item \textbf{Free end:} uniform dead traction $-P_\text{ref}/A$ in $x$
  direction on the face $x=L$. There are no further boundary conditions.
\item \textbf{Elements and mesh:} \texttt{SOLID HEX20} or \texttt{HEX27} with
  at least $2\times3$ elements in the cross-section and at least 50 along
  $L$. Alternatively \texttt{HEX8} with EAS, $4\times6\times100$.
\end{itemize}

\section{Analytical solution}

\paragraph{Euler--Bernoulli.}
From $EI\,v'''' + P\,v'' = 0$ with $v(0)=v'(0)=0$ and $M(L)=Q(L)=0$:
\begin{equation}
  P_{E,n} = \frac{(2n-1)^2\pi^2 EI}{4L^2},\qquad
  v_n(x) = 1-\cos\frac{(2n-1)\pi x}{2L},\qquad n=1,2,\dots
\end{equation}
This holds separately for both bending axes:
\begin{itemize}
\item weak axis, deflection in $z$: $I_y = b h^3/12 = \SI{2}{mm^4}$;
\item strong axis, deflection in $y$: $I_z = h b^3/12 = \SI{4.5}{mm^4}$.
\end{itemize}
The ratio $I_z/I_y = 2.25$ interleaves the two families, so the mode order
tests that modes are ordered correctly.
Torsional buckling of a solid rectangular section is far above the listed
range.

\paragraph{Shear correction (Engesser).}
\begin{equation}
  P_{T,n} = \frac{P_{E,n}}{1+P_{E,n}/(\kappa G A)},\qquad
  \kappa = \frac{10(1+\nu)}{12+11\nu} = 0.8497 \ \text{(Cowper)} .
\end{equation}
The Haringx variant differs from Engesser by less than \SI{0.02}{\percent}
for all listed modes. $P_{T,n}$ is the primary reference.

\section{Additional check: preloaded base state}
\begin{enumerate}
\item Run the nonlinear static preload with $P_\text{pre} = 0.5\,P_{T,1}$
  (e.g.\ \texttt{NUMSTEP: 5}).
\item Then run the buckling step with the same $P_\text{ref}$.
\end{enumerate}
The pre-buckling path of this structure is practically linear
($\varepsilon_\text{pre}<10^{-4}$), so the result must satisfy
$P_\text{pre} + \lambda_1^{\text{(pre)}}P_\text{ref} = \lambda_1^{(0)}P_\text{ref}$
up to $\approx 10^{-3}$ relative.

%======================================================================
\chapter{Case C -- Published benchmark: FEniCSx 3D solid}
\label{ch:fenics}
%======================================================================

\section{Source}
J.~Bleyer, M.~Pierre, \emph{Linear Buckling Analysis of a 3D solid},
Numerical Tours of Computational Mechanics with FEniCSx,
\url{https://bleyerj.github.io/comet-fenicsx/tours/eigenvalue_problems/buckling_3d_solid/buckling_3d_solid.html}.
The source is in \url{https://github.com/bleyerj/comet-fenicsx},
directory \texttt{tours/eigenvalue\_problems/buckling\_3d\_solid},
commit \texttt{5d06010c} (4 March 2025).

The tour solves exactly the classical problem $\K\bm{\phi} = -\lambda\K_\sigma\bm{\phi}$
with 3D solid elements, where $\K$ is the linear stiffness and $\K_\sigma$ is
built from the linear pre-stress $\bm{\sigma}_0$.

\section{Problem definition (identical to the benchmark)}
\begin{center}
\begin{tabular}{@{}ll@{}}
\toprule
Domain & box $[0,L]\times[0,b]\times[0,h]$, $L=1$, $b=0.01$ (in $y$), $h=0.03$ (in $z$) \\
Material & isotropic linear elastic, $E=10^3$, $\nu=0.3$ \\
Mesh & $50\times5\times5$ hexahedra, 2nd-order Lagrange ($Q_2$, 27 nodes) $\Rightarrow$ \texttt{SOLID HEX27} \\
DOFs & 36\,663 (including constrained DOFs) \\
Left end $x=0$ & fully clamped: $u_x=u_y=u_z=0$ on all face nodes \\
Right end $x=L$ & $u_y=u_z=0$ on all face nodes, $u_x$ free \\
Load & dead traction $\bm{t}=(-N_0,0,0)$ on the face $x=L$, $N_0=1$ (a stress) \\
Output & eigenvalue $\lambda$ = critical \emph{axial stress}; critical force $F_\text{cr}=\lambda\,b\,h$ \\
\bottomrule
\end{tabular}
\end{center}

\paragraph{Replicating it in 4C.}
\begin{itemize}
\item Generate the same structured $50\times5\times5$ \texttt{HEX27} mesh.
\item Boundary conditions:
  \begin{itemize}
  \item \texttt{DESIGN SURF DIRICH} on $x=0$ with \texttt{ONOFF [1,1,1]};
  \item \texttt{DESIGN SURF DIRICH} on $x=L$ with \texttt{ONOFF [0,1,1]};
  \item \texttt{DESIGN SURF NEUMANN} on $x=L$ with \texttt{VAL [-1,0,0]}.
  \end{itemize}
\item Material \texttt{MAT\_Struct\_StVenantKirchhoff} with $E=1000$,
  $\nu=0.3$, \texttt{KINEM nonlinear}, base state $\bu_0=\bm{0}$.
\end{itemize}
At $\bu_0=\bm{0}$ the 4C tangent $\KT(\bm{0})$ equals the benchmark's $\K$.

\section{Analytical solution (beam theory, as in the tour)}
Clamped--pinned Euler column: $EI\,v''''+F v''=0$ with $v(0)=v'(0)=0$ and
$v(L)=v''(L)=0$. This leads to the transcendental equation
\begin{equation}
  \tan\alpha_n = \alpha_n,\qquad F_{\text{cr},n} = \alpha_n^2\,\frac{EI}{L^2},
  \qquad \lambda_n = \frac{F_{\text{cr},n}}{b\,h} .
\end{equation}
The exact roots are
$\alpha_n/\pi = 1.430297,\ 2.459024,\ 3.470890,\ 4.477409,\ 5.481537,\ 6.484387$.
The two bending axes have
\begin{itemize}
\item weak axis, deflection in $y$: $I = h b^3/12$;
\item strong axis, deflection in $z$: $I = b h^3/12 = 9\times$ weak.
\end{itemize}
The shear-corrected (Engesser) values are listed as well.

\paragraph{Two remarks on the published tour.}
\begin{enumerate}
\item The tour quotes $\alpha_2\approx2.4509\pi$. The correct root of
  $\tan\alpha=\alpha$ is $\alpha_2 = 2.459024\pi$.
\item The tour compares FE mode $i$ with the $i$-th weak-axis beam mode for
  all $i$. FE modes 4 and 8, however, are strong-axis modes (dominant $u_z$),
  so from mode 4 on the tour's comparison pairs the wrong modes. The table in
  Chapter~\ref{ch:results} assigns each mode to its beam counterpart.
\end{enumerate}

\section{Numerical reference values}
The FE reference values were produced by running the unmodified benchmark
script. The only changes were an API update for dolfinx~0.9 (\texttt{.vector}
$\to$ \texttt{.x.petsc\_vec}) and additional output. The run used the
official container \texttt{dolfinx/dolfinx:v0.9.0} with SLEPc Krylov--Schur,
shift-invert $\sigma=0.01$, tolerance $10^{-12}$.
\begin{itemize}
\item The published mesh $50\times5\times5$ (36\,663 DOFs) is the primary
  reference for the code-to-code comparison.
\item The refined mesh $80\times8\times8$ (139\,587 DOFs) is the converged
  reference.
\item A $100\times5\times5$ run (72\,963 DOFs) differs from $80\times8\times8$
  by at most $4\times10^{-4}$ relative for modes 1--8, which indicates the
  remaining discretization error.
\item A $100\times10\times10$ run exceeded the 16\,GB of memory available and
  was not used.
\end{itemize}
The FE eigenvalues are not monotone under refinement, and mode~1 slightly
increases. This is expected, because the discrete pre-stress enters
$\K_\sigma$, so conforming FE eigenvalues are not bounds here. All eight
modes have at least \SI{99.8}{\percent} of their norm in the dominant
component ($u_y$ weak axis, $u_z$ strong axis).

%======================================================================
\chapter{Case D -- Periodic square-lattice cell (plausibility)}
\label{ch:periodic}
%======================================================================

\section{Purpose}
There is no clean analytical solution for a buckling problem that 4C's
periodic RVE conditions can represent: they require periodicity on all six
faces of a cuboid cell. This case is therefore checked by
\begin{enumerate}
\item exact identities that any correct implementation must satisfy, and
\item a closed-form frame-theory estimate as a plausibility band.
\end{enumerate}
It is the only case that exercises periodic boundary conditions and a
Dirichlet (prescribed macro-strain) perturbation load.

\section{Problem definition}
\begin{itemize}
\item \textbf{Unit cell:} $[0,L]\times[0,L]\times[0,d]$ with $L=\SI{20}{mm}$,
  $d=\SI{1}{mm}$, and a square hole $[w/2,\,L-w/2]^2$ through the full depth,
  $w=\SI{1}{mm}$. The material thus forms a square frame whose struts of total
  thickness $w$ are split in two halves across the periodic faces. Tiling
  gives an infinite square lattice with strut spacing $L$ and strut
  thickness $w$.
\item \textbf{Material:} St.\ Venant--Kirchhoff, $E=\SI{210000}{MPa}$,
  $\nu=0.3$.
\item \textbf{Periodic BCs:} \texttt{DESIGN SURF PERIODIC RVE 3D} \texttt{BOUNDARY CONDITIONS} on all six faces, penalty enforcement
  (\texttt{CONSTRAINT\_ENFORCEMENT: penalty}), automatic corner detection.
\item \textbf{Macro loading via the corner reference nodes} (N1 at the
  origin; N2, N4, N5 the corners in $+x$, $+y$, $+z$):
  \begin{itemize}
  \item N1: all components fixed;
  \item N2: $u_x = -\varepsilon_\text{ref}L$ with $\varepsilon_\text{ref}=10^{-3}$
    through a time function, $u_y=u_z=0$;
  \item N4: $u_x=u_z=0$, $u_y$ free;
  \item N5: $u_x=u_y=0$, $u_z$ free.
  \end{itemize}
  This gives uniaxial macro compression with free lateral contraction and no
  macro shear or rotation.
\item \textbf{Two model sizes:}
  \begin{itemize}
  \item D1: the $1\times1\times1$ cell;
  \item D2: a $2\times2\times1$ super-cell with the identical mesh in every
    sub-cell.
  \end{itemize}
\item \textbf{Mesh:} structured, \texttt{HEX20}/\texttt{HEX27}, at least
  2 elements across each half-strut, at least 40 along a strut, 1 element in
  $z$.
\item \textbf{Serial run:} 4C's RVE conditions are not implemented in
  parallel.
\end{itemize}

\section{Frame-theory estimate}
The struts in $x$ carry the whole load with strain $\varepsilon$. The struts
in $y$ are unloaded. Periodicity in $z$ with free uniform $\varepsilon_{zz}$
makes bending effectively plane strain, i.e.\ the bending stiffness is
$E w^3/[12(1-\nu^2)]$ per unit depth. The joints are treated as rigid blocks
of size $w$, so the free strut length is $\ell=L-w$.
\begin{itemize}
\item \textbf{D1 (period 1):} all joints rotate equally. The lowest mode is
  strut buckling between non-rotating joints (clamped--clamped, no sway):
  \begin{equation}
    \varepsilon_\text{cr}^{D1} = \frac{4\pi^2 w^2}{12(1-\nu^2)\,\ell^2} .
  \end{equation}
\item \textbf{D2 (period 2):} alternating transverse sway of neighbouring
  joint columns becomes admissible (clamped--guided):
  \begin{equation}
    \varepsilon_\text{cr}^{D2} = \frac{\pi^2 w^2}{12(1-\nu^2)\,\ell^2}
    = \tfrac14\,\varepsilon_\text{cr}^{D1} .
  \end{equation}
\end{itemize}
A periodic Euler--Bernoulli frame FE model (rigid joint zones, 16 elements
per strut) reproduces both expressions to four significant digits. Its spectrum also
contains every D1 eigenvalue within the D2 spectrum, which demonstrates
identity (I1) below. Joint flexibility, shear and 3D effects are not
included, so this estimate is only a plausibility band.

\section{Exact identities (independent of any reference value)}
\begin{enumerate}[label=(I\arabic*)]
\item \textbf{Super-cell containment:} every eigenvalue of D1 appears in the
  spectrum of D2, possibly with higher multiplicity. The reason is that every
  1-periodic field is also 2-periodic. Request at least $4\times$ as many
  modes for D2.
\item \textbf{Monotonicity:} $\lambda_1^{D2}\le\lambda_1^{D1}$.
\item \textbf{Symmetry:} loading in $y$ instead of $x$ (swap the roles of N2
  and N4) gives the same spectrum.
\item \textbf{Periodicity of modes:} $\bm{\phi}(\bm{x}^+)-\bm{\phi}(\bm{x}^-)$
  on paired nodes is of the order of the penalty error. It decreases by about
  $10\times$ when the penalty parameter is increased $10\times$.
\item \textbf{Linearity in the load:} $\lambda(2\varepsilon_\text{ref}) =
  \lambda(\varepsilon_\text{ref})/2$ and $\lambda(-\varepsilon_\text{ref})$
  flips the sign of the whole spectrum.
\item \textbf{Penalty convergence:} $\lambda_1(\kappa)$ converges as
  $\kappa\to\infty$. The change from $\kappa$ to $10\kappa$ must be at least
  $5\times$ smaller than the change from $\kappa/10$ to $\kappa$.
\end{enumerate}

%======================================================================
\chapter{Reference results to match}
\label{ch:results}
%======================================================================

This chapter collects all numbers and acceptance criteria.
Relative deviation: $\delta_i = |\lambda_i^\text{4C}-\lambda_i^\text{ref}|/\lambda_i^\text{ref}$.
The tolerances apply to the recommended meshes and are to be tightened into
regression values once the 4C results have been verified.

\section{Case A -- simply supported plate}

\subsection*{A1: $a=b=\SI{100}{mm}$, $t=\SI{1}{mm}$, $D=\SI{19230.769}{Nmm}$, $\pi^2D/b^2=\SI{18.980008}{N/mm}$}
\begin{center}
\begin{tabular}{@{}ccrrrr@{}}
\toprule
Mode & $(m,n)$ & $k_{mn}$ & Kirchhoff $\lambda$ & \refval{Mindlin $\lambda$ (reference)} & Mindlin/Kirchhoff \\
\midrule
1 & (1,1) & 4.000000 & 75.920034 & \refval{75.877241} & $-0.056\%$ \\
2 & (2,1) & 6.250000 & 118.625053 & \refval{118.458034} & $-0.141\%$ \\
3 & (3,1) & 11.111111 & 210.888983 & \refval{210.295972} & $-0.281\%$ \\
4 & (2,2) & 16.000000 & 303.680135 & \refval{302.996602} & $-0.225\%$ \\
5 & (4,1) & 18.062500 & 342.826403 & \refval{341.190800} & $-0.477\%$ \\
6 & (3,2) & 18.777778 & 356.402381 & \refval{355.100634} & $-0.365\%$ \\
\bottomrule
\end{tabular}
\end{center}
$\lambda$ = critical membrane force $N_x$ in \si{N/mm} ($N_\text{ref}=\SI{1}{N/mm}$),
equivalently the critical stress in \si{MPa} for $t=\SI{1}{mm}$.
Mode 7 is degenerate: $(1,2)$ and $(4,2)$ both have $k=25$. Do not compare
modes beyond 6.

\subsection*{A2: $a=\SI{300}{mm}$, $b=\SI{100}{mm}$, $t=\SI{1}{mm}$}
\begin{center}
\begin{tabular}{@{}ccrrrr@{}}
\toprule
Mode & $(m,n)$ & $k_{mn}$ & Kirchhoff $\lambda$ & \refval{Mindlin $\lambda$ (reference)} & Mindlin/Kirchhoff \\
\midrule
1 & (3,1) & 4.000000 & 75.920034 & \refval{75.877241} & $-0.056\%$ \\
2 & (4,1) & 4.340278 & 82.378509 & \refval{82.314032} & $-0.078\%$ \\
3 & (2,1) & 4.694444 & 89.100595 & \refval{89.064318} & $-0.041\%$ \\
4 & (5,1) & 5.137778 & 97.515066 & \refval{97.411294} & $-0.106\%$ \\
5 & (6,1) & 6.250000 & 118.625053 & \refval{118.458034} & $-0.141\%$ \\
6 & (7,1) & 7.628118 & 144.781743 & \refval{144.519114} & $-0.181\%$ \\
7 & (8,1) & 9.251736 & 175.598030 & \refval{175.197311} & $-0.228\%$ \\
8 & (1,1) & 11.111111 & 210.888983 & \refval{210.822928} & $-0.031\%$ \\
\bottomrule
\end{tabular}
\end{center}

\subsection*{Acceptance criteria}
\begin{itemize}
\item A1 modes 1--3: $\delta_i\le\SI{1}{\percent}$; A1 modes 4--6:
  $\delta_i\le\SI{2}{\percent}$.
\item A2 modes 1--8: $\delta_i\le\SI{1.5}{\percent}$, and the order of $(m,n)$
  must be exactly as listed. The spacing between modes 1 and 2 is only
  \SI{8.5}{\percent}, so a wrong order indicates an error.
\item Mode shape: half-wave count $(m,n)$ as listed; MAC $\ge 0.98$ against
  $\sin(m\pi x/a)\sin(n\pi y/b)$, evaluated with $u_z$ at the mid-surface
  nodes.
\item Mesh convergence: $\lambda_1$ decreases monotonically towards the
  Mindlin value under uniform refinement.
\item Load reversal: reversing the traction gives $\lambda_i\to-\lambda_i$
  (relative $10^{-8}$).
  Doubling the traction halves $\lambda_i$.
\end{itemize}

\section{Case B -- cantilever column}
$L=\SI{100}{mm}$, $b\times h=\SI{3}{mm}\times\SI{2}{mm}$; $\lambda$ = critical force in \si{N} ($P_\text{ref}=\SI{1}{N}$).
\begin{center}
\begin{tabular}{@{}clcrrr@{}}
\toprule
Mode & Axis (deflection) & $n$ & Euler $P_E$ [\si{N}] & \refval{Timoshenko $P_T$ [\si{N}] (reference)} & $P_T/P_E$ \\
\midrule
1 & weak ($z$) & 1 & 103.630846 & \refval{103.604771} & $-0.025\%$ \\
2 & strong ($y$) & 1 & 233.169404 & \refval{233.037442} & $-0.057\%$ \\
3 & weak ($z$) & 2 & 932.677616 & \refval{930.569806} & $-0.226\%$ \\
4 & strong ($y$) & 2 & 2098.524636 & \refval{2087.883909} & $-0.507\%$ \\
5 & weak ($z$) & 3 & 2590.771155 & \refval{2574.572274} & $-0.625\%$ \\
6 & weak ($z$) & 4 & 5077.911464 & \refval{5016.053142} & $-1.218\%$ \\
\bottomrule
\end{tabular}
\end{center}
\subsection*{Acceptance criteria}
\begin{itemize}
\item Modes 1--3: $\delta_i\le\SI{1}{\percent}$; modes 4--6:
  $\delta_i\le\SI{2}{\percent}$.
\item The order weak/strong/weak/strong/weak/weak must be exactly as listed.
  Dominant displacement component: $u_z$ for weak-axis modes, $u_y$ for
  strong-axis modes.
\item Mode shape: MAC $\ge 0.98$ against $1-\cos[(2n-1)\pi x/(2L)]$,
  evaluated along the centroidal axis.
\item Preload: with $P_\text{pre}=\SI{51.80}{N}$ ($\approx0.5P_{T,1}$),
  $P_\text{pre}+\lambda_1^{(\text{pre})}\cdot\SI{1}{N}$ matches
  $\lambda_1^{(0)}$ within $10^{-3}$ relative.
\end{itemize}

\section{Case C -- FEniCSx benchmark}
$\lambda$ = critical compressive traction (stress) on the end face; $N_0=1$.
The FEniCSx columns are the numerical reference; the beam columns are
analytical plausibility values.

\begin{center}
\small
\begin{tabular}{@{}cc c r r r r@{}}
\toprule
Mode & Axis & $n$ & \refval{FEniCSx $50{\times}5{\times}5$} & FEniCSx $80{\times}8{\times}8$ & Euler beam & Timoshenko beam \\
\midrule
1 & weak ($y$) & 1 & \refval{0.167963169} & 0.168078306 & 0.16825607 & 0.16816949 \\
2 & weak ($y$) & 2 & \refval{0.496955439} & 0.496866498 & 0.49732930 & 0.49657360 \\
3 & weak ($y$) & 3 & \refval{0.987888956} & 0.987687419 & 0.99083224 & 0.98783717 \\
4 & strong ($z$) & 1 & \refval{1.500093336} & 1.500692451 & 1.51430464 & 1.50732006 \\
5 & weak ($y$) & 4 & \refval{1.642490988} & 1.640922961 & 1.64881509 & 1.64053797 \\
6 & weak ($y$) & 5 & \refval{2.455332621} & 2.451987909 & 2.47128677 & 2.45273882 \\
7 & weak ($y$) & 6 & \refval{3.429242204} & 3.421467271 & 3.45824987 & 3.42203704 \\
8 & strong ($z$) & 2 & \refval{4.392907156} & 4.392598990 & 4.47596370 & 4.41548720 \\
\bottomrule
\end{tabular}
\end{center}

\subsection*{Acceptance criteria}
\begin{itemize}
\item \textbf{Code-to-code (primary):} 4C with \texttt{HEX27},
  $50\times5\times5$, identical BCs, against FEniCSx $50\times5\times5$:
  $\delta_i\le10^{-4}$ for modes 1--8. The only allowed difference is the
  quadrature rule.
\item Mode identification (dominant component) exactly as listed. The
  dominant component carries at least \SI{99}{\percent} of the mode norm.
\item 4C $80\times8\times8$ \texttt{HEX27} against FEniCSx $80\times8\times8$: same tolerances as above. Any finer 4C mesh against FEniCSx $80\times8\times8$: $\delta_i\le10^{-3}$.
\end{itemize}

\section{Case D -- periodic lattice cell}
$L=\SI{20}{mm}$, $w=\SI{1}{mm}$, $\ell=L-w=\SI{19}{mm}$, $\varepsilon_\text{ref}=10^{-3}$; $\lambda=\varepsilon_\text{cr}/\varepsilon_\text{ref}$.
\begin{center}
\begin{tabular}{@{}llcc@{}}
\toprule
Model & Mode type & \multicolumn{2}{c}{Frame estimate $\lambda$} \\
 & & point joints, $\ell=L$ & \refval{rigid joints, $\ell=L-w$} \\
\midrule
D1 ($1\times1\times1$) & clamped--clamped, no sway & 9.0381 & \refval{10.0145} \\
D2 ($2\times2\times1$) & alternating sway & 2.2595 & \refval{2.5036} \\
\bottomrule
\end{tabular}
\end{center}
\subsection*{Acceptance criteria}
\begin{itemize}
\item \textbf{Identities (hard):}
  \begin{itemize}
  \item (I1) every D1 eigenvalue, up to the requested count, is found in the
    D2 spectrum within $10^{-6}$ relative (eigensolver tolerance);
  \item (I2) $\lambda_1^{D2}<\lambda_1^{D1}$;
  \item (I3) $x$- and $y$-loading spectra agree within $10^{-6}$;
  \item (I4)--(I6) as defined in Chapter~\ref{ch:periodic}.
  \end{itemize}
\item \textbf{Plausibility (soft):}
  \begin{itemize}
  \item $\lambda_1^{D1}$ and $\lambda_1^{D2}$ within $\pm\SI{15}{\percent}$ of
    the rigid-joint frame estimate;
  \item $\lambda_1^{D2}/\lambda_1^{D1}$ between 0.2 and 0.3 (frame theory:
    0.25);
  \item D1 mode: $x$-struts buckle with the cell's joints non-rotating;
    D2 mode: neighbouring joint columns sway in opposite $y$ directions.
  \end{itemize}
\end{itemize}

\end{document}
